% 页面
\ifdefstring{\Ysy@type}{universal}{
	\geometry{
		a4paper,
		top=25.4mm, bottom=25.4mm,
		left=20mm, right=20mm,
		headheight=2.17cm,
		headsep=4mm,
		footskip=12mm
	}
}{\relax}
\ifdefstring{\Ysy@type}{paper}{
	\geometry{
		a4paper,
		top=25.4mm, bottom=25.4mm,
		left=20mm, right=20mm,
		headheight=2.17cm,
		headsep=4mm,
		footskip=12mm
	}
}{\relax}
\ifdefstring{\Ysy@type}{flab}{
	\geometry{top=28mm,bottom=28mm,left=15mm,right=15mm}
}{\relax}


% 超链接
\hypersetup{
	breaklinks,
	unicode,
	linktoc=all,
	bookmarksnumbered=true,
	bookmarksopen=true,
	pdfkeywords={YsyOneForAll},
	colorlinks,
	linkcolor=LinkColor,
	citecolor=LinkColor,
	urlcolor=LinkColor,
	plainpages=false,
	pdfstartview=FitH,
	pdfborder={0 0 0},
	linktocpage
}

% 引用
\crefformat{figure}{#2{\textcolor{EmphColor}{\textbf{\YsyFigureName#1}}}#3} % 图片的引用格式
\crefformat{equation}{#2{\textcolor{EmphColor}{(#1)}}#3} % 公式的引用格式
\crefformat{table}{#2{\textcolor{EmphColor}{\textbf{\YsyTableName#1}}}#3} % 表格的引用格式
\crefformat{enumi}{#2{\textcolor{EmphColor}{\textbf{[#1]}}}#3} % 文献的引用格式

% 图片、表格名称
\captionsetup[figure]{labelfont={color=StructureColor!70!black,bf},name=\YsyFigureName} % 图片形式
\captionsetup[table]{labelfont={color=StructureColor!70!black,bf},name=\YsyTableName} % 表格形式

% 页眉页脚
\ifdefstring{\Ysy@type}{universal}{
	\pagestyle{fancy}
	\fancyhf{} 
	\renewcommand{\headrulewidth}{0pt}
	\renewcommand{\footrule}{\color{StructureColor}\hrule width\textwidth height 2pt} 
	\fancyhead[R]{\textcolor{StructureColor!70!LinkColor}{\textbf{\TitleName}}} 
	\fancyfoot[L]{\textcolor{StructureColor!70!LinkColor}{\textbf{\AuthorName}}}
	\fancyfoot[R]{\textcolor{StructureColor!70!LinkColor}{\textbf{\thepage}}}
}{\relax}
\ifdefstring{\Ysy@type}{flab}{
	%	页眉页脚设置
	\usepackage{fancyhdr}
	\fancypagestyle{plain}{\pagestyle{fancy}}
	\pagestyle{fancy}
	\lhead{\kaishu 中山大学物理与天文学院基础物理实验} % 左边页眉，学院 + 课程
	\rhead{\kaishu \TitleName} % 右边页眉，实验报告标题
	\cfoot{\thepage} % 页脚，中间添加页码
	\setlength{\headheight}{13.6pt}
}{\relax}
\ifdefstring{\Ysy@type}{paper}{
	%	页眉页脚设置
	\pagestyle{fancy}
	% 清除原有样式
	\fancyhf{}
	% 页眉设置
	\renewcommand{\headrulewidth}{0.4pt}  % 双线第一条
	\renewcommand{\headrule}{
		{\color{StructureColor}\hrule height 0.4pt}
		{\vskip1pt}
		{\color{StructureColor}\hrule height 0.4pt}
		{\vskip2pt}
	}
	\fancyhead[C]{%
		\textcolor{StructureColor!70!LinkColor}{\textbf{\normalsize\TitleName} \\[1pt]
		{\small\AddInfo}}
	}
	% 页脚设置
	\fancyfoot[C]{\textcolor{StructureColor!70!LinkColor}{-- \textbf{\thepage} --}}
	\renewcommand{\footrulewidth}{0pt} % 无脚线
}{\relax}

% 列表
% 有序列表
\setlist[enumerate,1]{label=\textcolor{StructureColor!80!black}{\arabic*.}, font=\bfseries\color{StructureColor!80!black}} 
\setlist[enumerate,2]{label=\textcolor{StructureColor!80!black}{(\arabic*)}, font=\bfseries\color{StructureColor!80!black}}
\setlist[enumerate,3]{label=\textcolor{StructureColor!80!black}{\roman*.}, font=\bfseries\color{StructureColor!80!black}}
\setlist[enumerate,4]{label=\textcolor{StructureColor!80!black}{(\roman*)}, font=\bfseries\color{StructureColor!80!black}}
% 无序列表
\ifYsy@fancyitem
\setlist[itemize,1]{label=\liuyin}
\else
\setlist[itemize,1]{label={\eitemi}}
\fi
\setlist[itemize,2]{label={\eitemii}}
\setlist[itemize,3]{label={\eitemiii}}

% 章节
\ifdefstring{\Ysy@type}{universal}{
	\ifYsy@fancysection
	\titleformat{\section}[block]
	{\normalfont\bfseries\color{StructureColor!80!black}\huge}
	{\liuyin[1.5ex]}
	{0.618em}
	{}
	[
	\vspace{-0.8cm} % 负值缩小间距
	\noindent\textcolor{StructureColor}{\rule{0.382\textwidth}{2pt} }
	\vspace{7pt}
	]
	\else
	\titleformat{\section}[block]
	{\normalfont\bfseries\color{StructureColor!80!black}\huge}
	{\thesection}
	{0.618em}
	{}
	\fi
}{\relax}
\ifdefstring{\Ysy@type}{flab}{
	\titleformat{\section}[block]
	{\centering\normalfont\bfseries\color{StructureColor!80!black}\huge}{}{1em}{}
}{\relax}
\titleformat{\subsection}[block]
{\normalfont\bfseries\color{StructureColor!80!black}\LARGE} % 设置字体、大小、颜色
{\thesubsection} % 设置subsection编号的格式
{0.618em} % 编号和标题之间的间距
{}
\setcounter{secnumdepth}{3}
\titleformat{\subsubsection}[block]
{\normalfont\bfseries\color{StructureColor!80!black}\Large}
{\thesubsubsection}
{0.618em}
{}
\ifdefstring{\Ysy@type}{paper}{
	\titleformat{\section}[block]
	{\normalfont\bfseries\color{StructureColor!80!black}\huge}
	{\thesection}
	{0.618em}
	{}
}{\relax}

% 目录
\renewcommand\tableofcontents{%
	\titleformat{\section}[block]
	{\normalfont\bfseries\color{StructureColor!80!black}\huge}
	{}
	{0.618em}
	{}
	
	\hypersetup{linktoc=all, linkcolor=black}
	\thispagestyle{empty}
	\section*{\centerline{\YsyContentsName}}
	\@starttoc{toc}
	\hypersetup{linkcolor=LinkColor}
	\ifdefstring{\Ysy@type}{universal}{
		\ifYsy@fancysection
		\titleformat{\section}[block]
		{\normalfont\bfseries\color{StructureColor!80!black}\huge}
		{\liuyin[1.5ex]}
		{0.618em}
		{}
		[
		\vspace{-0.8cm} % 负值缩小间距
		\noindent\textcolor{StructureColor}{\rule{0.382\textwidth}{2pt} }
		\vspace{7pt}
		]
		\else
		\titleformat{\section}[block]
		{\normalfont\bfseries\color{StructureColor!80!black}\huge}
		{\thesection}
		{0.618em}
		{}
		\fi
	}{\relax}
	\ifdefstring{\Ysy@type}{flab}{
		\titleformat{\section}[block]
		{\centering\normalfont\bfseries\color{StructureColor!80!black}\huge}{}{1em}{}
	}{\relax}
	\ifdefstring{\Ysy@type}{paper}{
		\titleformat{\section}[block]
		{\normalfont\bfseries\color{StructureColor!80!black}\huge}
		{\thesection}
		{0.618em}
		{}
	}{\relax}
}

% 图片路径
\graphicspath{{./images/}{./ysy-latex/theme}}

% 参考文献
\newcommand{\YsyRef}[7]{
	\item\label{ref:#1} #2\quad#3[#4]. \emph{#5}, #6, #7.
}

\newcommand{\YsyReferencesPage}[1]{%
	\titleformat{\section}[block]
	{\normalfont\bfseries\color{StructureColor!80!black}\huge}
	{}
	{0.618em}
	{}
	
	\section*{\centerline{\YsyReferencesName}}
	\addcontentsline{toc}{section}{\YsyReferencesName}
	\noindent
	\begin{enumerate}[label=\textbf{[\arabic*]}, ref=\arabic*]
		#1
	\end{enumerate}
	
	\ifdefstring{\Ysy@type}{universal}{
		\ifYsy@fancysection
		\titleformat{\section}[block]
		{\normalfont\bfseries\color{StructureColor!80!black}\huge}
		{\liuyin[1.5ex]}
		{0.618em}
		{}
		[
		\vspace{-0.8cm} % 负值缩小间距
		\noindent\textcolor{StructureColor}{\rule{0.382\textwidth}{2pt} }
		\vspace{7pt}
		]
		\else
		\titleformat{\section}[block]
		{\normalfont\bfseries\color{StructureColor!80!black}\huge}
		{\thesection}
		{0.618em}
		{}
		\fi
	}{\relax}
	\ifdefstring{\Ysy@type}{flab}{
		\titleformat{\section}[block]
		{\centering\normalfont\bfseries\color{StructureColor!80!black}\huge}{}{1em}{}
	}{\relax}
	\ifdefstring{\Ysy@type}{paper}{
		\titleformat{\section}[block]
		{\normalfont\bfseries\color{StructureColor!80!black}\huge}
		{\thesection}
		{0.618em}
		{}
	}{\relax}
}

% 附录
\newcommand{\YsyAppendixPage}[1]{%
	\titleformat{\section}[block]
	{\normalfont\bfseries\color{StructureColor!80!black}\huge}
	{}
	{0.618em}
	{}
	
	\renewcommand\thesubsection{\Alph{subsection}}
	\titleformat{\subsection}[block]
	{\normalfont\bfseries\color{StructureColor!80!black}\LARGE} % 设置字体、大小、颜色
	{\thesubsection} % 设置subsection编号的格式
	{0.618em} % 编号和标题之间的间距
	{}
	
	\section*{\centerline{\YsyAppendixName}}
	\addcontentsline{toc}{section}{\YsyAppendixName}
	
	\setcounter{subsection}{0}
	#1
	
	\ifdefstring{\Ysy@type}{universal}{
		\ifYsy@fancysection
		\titleformat{\section}[block]
		{\normalfont\bfseries\color{StructureColor!80!black}\huge}
		{\liuyin[1.5ex]}
		{0.618em}
		{}
		[
		\vspace{-0.8cm} % 负值缩小间距
		\noindent\textcolor{StructureColor}{\rule{0.382\textwidth}{2pt} }
		\vspace{7pt}
		]
		\else
		\titleformat{\section}[block]
		{\normalfont\bfseries\color{StructureColor!80!black}\huge}
		{\thesection}
		{0.618em}
		{}
		\fi
	}{\relax}
	\ifdefstring{\Ysy@type}{flab}{
		\titleformat{\section}[block]
		{\centering\normalfont\bfseries\color{StructureColor!80!black}\huge}{}{1em}{}
	}{\relax}
	\ifdefstring{\Ysy@type}{paper}{
		\titleformat{\section}[block]
		{\normalfont\bfseries\color{StructureColor!80!black}\huge}
		{\thesection}
		{0.618em}
		{}
	}{\relax}
	
	\renewcommand\thesubsection{\thesection.\arabic{subsection}}
	\titleformat{\subsection}[block]
	{\normalfont\bfseries\color{StructureColor!80!black}\LARGE} % 设置字体、大小、颜色
	{\thesubsection} % 设置subsection编号的格式
	{0.618em} % 编号和标题之间的间距
	{}
}